\documentclass{article}
\usepackage{amsmath, amssymb, amsthm}
\usepackage{hyperref}
\usepackage{geometry}
\geometry{margin=1in}

\title{Erd\H{o}s Problem \#1135}
\date{Last updated: 2026-01-11}
\author{Source: erdosproblems.com}

\begin{document}
\maketitle

\noindent\textbf{Status:} OPEN \quad \textbf{Prize: \$500}

\noindent\textbf{Tags:} number theory

\medskip
\noindent\textbf{Note:} Collatz conjecture

\medskip

\noindent\textbf{Problem Statement:}

Define $f:\mathbb{N}\to \mathbb{N}$ by $f(n)=n/2$ if $n$ is even and $f(n)=\frac{3n+1}{2}$ if $n$ is odd. Given any integer $m\geq 1$ does there exist $k\geq 1$ such that $f^{(k)}(m)=1$?




The infamous Collatz conjecture . For a detailed discussion of the history and theory surrounding this problem we refer to the overview by Lagarias \cite{La10}. This is not a problem due to Erd\H{o}s; it was first devised by Collatz before 1952. Erd\H{o}s referred to this problem on several occasions as 'hopeless'. As Lagarias \cite{La16} notes, the closest Erd\H{o}s ever came to working on problems of this nature is the theorem described in the remarks to [1134] . It is often claimed that Erd\H{o}s offered \$500 for a solution to this problem; this claim originated in a survey article by Lagarias \cite{La85}. Lagarias reported, in personal communication, that this came from a conversation he had with Erd\H{o}s and Graham around 1983, in which Graham asked Erd\H{o}s to make an estimate of what value Erd\H{o}s would put the problem on his prize scale, to which Erd\H{o}s replied \$500. Therefore, strictly speaking, Erd\H{o}s never offered \$500 specifically as a prize, but we include this prize value here for comparing those problems which Erd\H{o}s rated as 'prize problems'. This is Problem E16 in Guy's collection \cite{Gu04}, in which Guy quotes Erd\H{o}s as saying "Mathematics may not be ready for such problems".




References

[Gu04] Guy, Richard K., Unsolved problems in number theory . (2004), xviii+437.

[La10] Lagarias, Jeffrey C., The {$3x+1$} problem: an overview . (2010), 3--29.

[La16] Lagarias, Jeffrey C., Erd\H{o}s, {K}larner, and the {$3x+1$} problem . Amer. Math. Monthly (2016), 753--776.

[La85] Lagarias, Jeffrey C., The {$3x+1$} problem and its generalizations . Amer. Math. Monthly (1985), 3--23.

\noindent\textbf{Related OEIS sequences:} A006370, A008908


\bigskip
\noindent\small{Source: \url{https://www.erdosproblems.com/1135}}
\end{document}
